\subsection{Radius, amplitude, modular, and five-channel closure}
\label{sec:bt3486-3499}

An exact basis-exchange certificate in the minimum-defect Cayley system produces a
247-term circuit with 246 nonzero basis coordinates.  Since the coefficient alphabet
has only $3^5-1=242$ nonzero vectors, a nonzero circuit shift cancels at least two
basis coefficients while adding one extra generator.  The covering-radius interval
therefore improves to
\[
  \boxed{389\leq D_{H^1}\leq435}.
\]
The exact radius remains open.

A five-channel amplitude-bearing magnetic family was then screened on a 63,869-point
dyadic grid.  Its best member has weights
\[
 \left(-\frac{21}{32},\frac{43}{32},-\frac{27}{32},-\frac{15}{16},1\right)
\]
and generalized Hoffman ratio approximately $8.9024605237$.  After scaling by 32,
its exact characteristic polynomial contains $(x+128)^{17}(x-64)^{21}$ and a
seven-dimensional residual whose roots all lie strictly in $(-128,1024)$.
Consequently the unscaled matrix has $\lambda_{\min}=-4$ and
$\lambda_{\max}<32$, proving that this winner remains strictly below nine.  This is
a finite-grid no-go, not an optimization of the unrestricted real cone.

Over $\mathbb F_3$, the 240 face boundaries form a totally isotropic subspace of the
480-dimensional flat edge space.  They are disjoint from the 44-dimensional
coboundary space, giving the exact filtration
\[
  0<R_{240}<Z^1_{480}<\mathbb F_3^{720},\qquad
  \dim Z^1/(R+B^1)=196.
\]
No modular irreducible name is assigned without a composition-series certificate.

The full 135-state quotient admits an executable Fourier organization
\[
  135=27\times5.
\]
For $c(0)=2$ and $c(1)=c(2)=-1$, the four barycentric channels have symbol
$c(k_1)K_x+c(k_2)K_y+c(k_3)I_4$, while the hidden channel is
$-c(k_1)-c(k_2)+c(k_3)$.  Exhausting all 675 symbol entries recovers
\[
  10^1,\;7^6,\;4^{22},\;1^{44},\;(-2)^{42},\;(-5)^{20}.
\]

The finite geometry itself forms a stabilizer chain
\[
  5\xrightarrow{\rm hinge}4\xrightarrow{\rm pairing}3.
\]
Choosing one of the five Clebsch/$D_5$ coordinate directions leaves the $S_4$
action on the four points of the $PG(2,3)$ null conic.  The three perfect matchings
of those four points carry the quotient $S_4/V_4\cong S_3$ and match the three
projective $A_2(\mathbb F_3)$ root directions and the far/middle/active Fano gauges.
This is a finite permutation correspondence, not a physical $\mathrm{Spin}(10)$
claim.

Finally, the fixed linear $[13,5,5]$ controller encoder has a unique minimum
coordinate-orbit completion of length 28.  The completed code is
\[
  \boxed{[28,5,11]},
\]
so fifteen added coordinates are necessary and sufficient for coordinate-permutation
$S_3$ equivariance of that encoder.  The Clebsch closed-neighborhood matrix is the
incidence matrix of a symmetric $2$-$(16,6,2)$ biplane, with Smith form
$1^6 2^4 4^5 12$ and the modular identity $B^2=I-J$ over $\mathbb F_3$.
Its zero and single-point syndromes have minimum distance six.
